\section{The principal row lattices and rational pairing on the selected family}

% BEGIN R88L1
\paragraph{R88L1. Rational row coordinates at each attained tuple on the chosen bundle.}
Use X87S1 with the actual joint selector, and identify the native
datum with the datum of the linked point. Its witness contains a
twisted input $x$ and the assertion $\mathrm{Holds}(x.R,x.F)$.
In particular it contains the primitive row identities, the
well-formedness input for the finite family, and separation of its
types. The input $x$ itself contains the character and model
identifications, the readings $Y,Z,C$, and the first division
comparison $x.\mathrm{div}$. The sharp and the base parts of this last
comparison give all three agreements for the forms, integral
generators and map $\iota$ on each of the two levels. These hypotheses
are present for every guarded member of the selected class. We spell
out the portion of their finite-dimensional consequence needed below.

At either level and each tuple $t$, let $V_t$ be the matrix of
the rational good part in the level coordinates, $P_t$ its
quotient map to the non-kernel coordinates, and $J_t$ the map
from the latter complex coordinates to the level target. Then
$P_t$ is onto, $J_t$ is injective and all $J_t$ have the same
image. Here the finite level projector and the good projector
identify the range of $V_t$ with the good level space. The
kernel cut and its orthogonal complement are rational: evaluating
the given separating rational tests on the rational basis of
the source expresses the kernel by equations over the coefficient
field $K$, and its orthogonal complement by the rational Hermitian
Gram matrix. Intersect with the level space and choose a basis of
the resulting quotient. $P_t$ is its coordinate projection,
and the corresponding columns lie in the range of the adjoint
of the row operator $\Theta_t$. This proves injectivity of $J_t$,
since $\ker\Theta_t\cap\operatorname{im}\Theta_t^*=0$.
The primitive equivariance and level identities put its image in
the prescribed common level target. In the other direction, this
target is the sum of its given isotypic parts. On each such
part the primitive row identity gives
$\Theta_t\Theta_t^* e=\kappa\prod_j\lambda_{j,t}e$ with positive
scalar. The indicated adjoint vector is in the good level space
by equivariance and the stored upper-level identity. Its coordinate
basis in the cut gives a preimage of $e$ after division by the
scalar. This proves equality of the ranges, also when a part has
dimension greater than one.

Let $t_0$ be the base tuple. To compute the matrix $M_t$ that
satisfies $J_{t_0}M_t=J_t$, apply the given separating row
functionals to both sides. At the embedding used for these
coordinates the resulting matrices have entries in $K$, by the
row test on the rational basis vectors. The matrix of tests on the
$t_0$ columns has full column rank, by separation and injectivity.
Choose an invertible square minor and solve for each column.
Thus the existing complex change of coordinates is the embedding
of $M_t$ over $K$. Its determinant is nonzero by the common-image
and injectivity assertions above, and $M_{t_0}=1$.
The equation in $x.\mathrm{div}$ uses exactly this algebraic
matrix in the original native embedding:
\[
 \iota_{Y(t)}\big(j_{A,t}(V_tz)\big)
       = j_{EE}\big(e_{RR}(M_tP_tz)\big)\qquad(z\text{ over } K).
\]
Here $j_{A,t}$ and $e_{RR}$ are the readings on the level in
question; at the base level they are the stored base versions.
This explicit equality incorporates any change between the model
embedding and the native one.
% END R88L1

% BEGIN R88L2
\paragraph{R88L2. Fullness of both original principal row lattices.}
Fix an attained native tuple $\mu_{\rm At}(w)$. Its recorded
coordinate tuple $t$ satisfies $Y(t)=\mu_{\rm At}(w)$.
Let $y$ be an arbitrary member of the native rational target.
By R88L1 and the isomorphism $e_{RR}$, choose rational $z$ with
$e_{RR}(M_tP_tz)=y$. In the level coordinate space the integral
form module is a full lattice. Indeed the underlying finite
class-function lattice is full by the finite-family hypotheses;
intersecting it with a rational level-character space still spans
that space, since one may multiply any rational vector by a
nonzero integer to clear its coordinates in the lattice. Transport
by the given level basis. Multiply $V_tz$ by a nonzero
$h\in\mathcal O_K$ to put it in this lattice. The integral
reading now places $j_{A,t}(hV_tz)$ in the $\mathcal O_K$-span
of the native distinguished integral generators. The reading $Y$
makes the actual carrier stable under every element of
$\mathcal O_K$ for this level. It contains its generators and is
closed under addition, hence it contains this vector. Also the
vector is in the native good space by the projector agreement.

For any good vector $v$, the difference between $v$ and its
projection to $Q$ lies in the kernel of $\iota$. Indeed within the
good space $Q$ is the orthogonal complement of its intersection with
that kernel. Applying the definition of $\mathrm{RRset}$ and R88L1
therefore gives $j_{EE}(hy)\in\mathrm{RRset}(\mu_{\rm At}(w))$.
The rational span of $A_{\mu_{\rm At}(w)}$ in B87H2 contains $y$.
This holds for every attained tuple at the sharp level and at the
base level. B87H1--2 now give $D.\mathrm{RRFull}$ and the original
$\mathrm{BaseRowLatticeFull}$ on this very member. In particular
the fifth entry of $\mathrm{LinkDIV}_{\rm at}$ holds simultaneously
at all indices for $x\ge\mathcal F.x_1+1$ on the selected class.
% END R88L2

% BEGIN R88L3
\paragraph{R88L3. The selected rational values equal the given map.}
At the sharp level use the base tuple $t_0$. For any rational
$y$ and $\phi$ choose $z$ and $v$ with $e_{RR}(P_{t_0}z)=y$ and
$e_{N,u}(v)=\phi$. The base tuple has $M_{t_0}=1$.
Write $J(z)=j_{A,t_0}(V_{t_0}z)$. The kernel of $\iota$ inside
the good space is the complex span of $J(\ker P_{t_0})$.
In fact $J$ gives the good space upon taking complex spans, and
R88L1 identifies its composite with $\iota$ with the matrix
$P_{t_0}$ followed by the native coefficient embedding in a
complex-independent basis. Scalar extension of a matrix kernel
therefore gives precisely the asserted kernel.

In this ambient space the native Gram pairing of rational input
columns has entries in the image of $K$ by the forms reading.
Choose a rational basis of independent $J$-columns for
$J(\ker P_{t_0})$. The Gram matrix is nonsingular:
it is the Gram matrix of complex-independent vectors for a positive
Hermitian form. Solving the Gram equations for the orthogonal
projection of $J(z)$ to this kernel therefore gives $J(z')$ for
$z'\in\ker P_{t_0}$ over $K$. The field is stable under the
conjugation used for this calculation. Hence $J(z-z')$ belongs to
$Q$, and $\iota J(z-z')=j_{EE}y$.
The native section in the definition of $\mathrm{Bu}$ sends this
last vector to $J(z-z')$, since it inverts $\iota$ on $Q$.
Using the period reading $C.\mathrm{pt}$ at $t_0$ gives
\[
 \mathrm{Bu}_u(j_{N,u}\phi,j_{EE}y)
   =\mathrm{perCol}_{u,\mu_0}(j_{N,u}\phi,J(z-z'))
   =\mathrm{emb}_K\bigl(\mathrm{Pt}_{C.wU(u),t_0}(v,z-z')\bigr).
\]
Thus the native values lie in the original coefficient field at
all rational arguments.

Finally the norm extension with the list, retained by the common
extension of X87S2, contains the norm extension with the law for
$Bk$. Project the norm-extension witness first to the extension with
the list, then to its extension with the law, and take its law
component. This is the implication, on the same $D$, from the
rational-value property just proved to $D.\mathrm{BuRational}$.
It identifies the above rational values with the recorded
$D.\mathrm{div}.Bk$ under the given embedding. This establishes
the rational half of the original Bessel entry on all the guarded
indices of the selected class. Its common-unitary clause is a
separate comparison of the two complex period pairings.
% END R88L3
